\section{Results Discussion}
The updated results change the story from a synthetic-only placeholder into a genuine prototype benchmark. The important result is not that a neural model already dominates, because it does not. The important result is that the repository now supports a full progression from synthetic supervision to weak-label bootstrapping and exposes where each branch succeeds or fails.

The new tri-modal emotion benchmark clarifies one point from the advisor discussion: for upstream emotion understanding, multimodal fusion really does matter. On the synthetic 12-way fine-emotion task, the full face+motion+audio encoder reaches 97.59\% accuracy and 97.58\% Macro-F1, while the three single-modality ablations drop sharply. That does not prove real-world perception is solved, because the benchmark is synthetic, but it is enough to justify treating the tri-modal encoder as a serious upstream component rather than a cosmetic add-on.

On the large synthetic benchmark, both classical and neural timing models become credible, and the current split is no longer the earlier degenerate case. The structured baseline reaches 85.56\% balanced accuracy and 79.11\% timing Macro-F1. The multitask temporal model reaches 86.88\% and 77.21\%, while also delivering the best state-prediction result at 0.069 MAE. The exploratory joint model is strongest overall on synthetic timing at 86.90\% balanced accuracy and 81.06\% Macro-F1. This is enough to show that the descriptor-sequence backbone is learning a meaningful three-class timing task, but it still does not justify a claim about raw-video understanding in the wild.

The mixed results are harder to interpret and more important scientifically. Aggregate combined metrics can look strong, especially for the joint model at 85.79\% balanced accuracy and 80.45\% timing Macro-F1, but those combined columns are driven mostly by the 2{,}400 synthetic test windows rather than the 62 weaklabel windows. Once we isolate the held-out weaklabel split, the bottleneck is clear: timing Macro-F1 is 36.29\% for the joint model, 27.29\% for the structured baseline, and only 7.62\% for the multitask branch. This gap is the most honest summary of the current system's maturity.

The joint model should therefore be interpreted as the strongest current offline branch, not as a solved timing-to-strategy system. It is the only branch that remains competitive on synthetic timing while also predicting strategy, and its combined strategy Macro-F1 reaches 82.66\%. Yet the held-out weaklabel strategy result is only 18.06\%, which means the utterance-planning path is still mostly validated by synthetic or synthetic-dominated labels. That is enough to justify continued work on the interface, but not enough to claim robust support generation.

The structured baseline remains important even though it is no longer the strongest overall model. Its weaklabel timing Macro-F1 of 27.29\% and weaklabel strategy Macro-F1 of 22.50\% show that hand-engineered policy features still transfer better than the multitask branch in some log-derived settings. This confirms the upper-level audit: the repository is now a solid prototype and a reproducible research pipeline, but it is still short of the standard required for top-tier claims about natural video understanding, validated personalization, or real-user proactive care.
